\documentclass[UTF8,fontset=none,11pt,a4paper]{ctexart}
\usepackage{ipara-handbook}
\handbooksetup{DS-B02 Index Strategy × Workload}{408 · 数据结构 Bridge B02}{Workload · Redundancy · Cost Vector}
\begin{document}
\handbooktitle{Index Strategy × Workload}{查询能力来自被维护的冗余，选择必须由操作分布决定}{408 · 数据结构 Internal Bridge B02}

\begin{corebox}{桥梁接口}
\[
\text{Workload}\to\text{Required Operations}\to\text{Maintained Order/Redundancy}\to\text{Search/Update/Space/I/O Cost}\to\text{Choice}.
\]
有序索引维护比较关系，Hash 维护槽位映射；二者不是“谁更快”的绝对排名，而是用不同冗余换取不同查询和更新能力。
\end{corebox}
\begin{methodbox}{Owners 与停止边界}
DS09 Owns 折半查找、BST/AVL/B/B+ 机制；DS10 Owns Hash 冲突与装填；Atlas Foundation Owns cost vector。本桥不重复这些机制，只定义比较维度与选型证据。
\end{methodbox}
\tableofcontents\newpage

\section{先把 workload 写成操作向量}
记录精确等值、范围、前驱后继、排序遍历、插入删除更新的频率，以及是否要求最坏保证、稳定顺序、内存上限或块 I/O。若只问等值存在性且更新频繁，Hash 的平均成本可能占优；若要范围扫描或有序输出，有序索引保留必要顺序。

\section{冗余的交换关系}
维护排序或平衡信息会增加插入/删除修复；维护 Hash 槽位会支付冲突、装填和重散列。总成本可写成
\[
C_{total}=\sum_i f_i C_{query,i}+\sum_j g_j C_{update,j}+C_{space}+C_{I/O}.
\]
这不是精确机器计时，而是防止只比较一个操作的决策框架。
\begin{examplebox}{同一数据的两种选择}
日志按 user\_id 精确查找且很少范围查询：Hash 避免维护全序；若要求按时间区间扫描并支持前驱/后继，则需要有序索引或额外排序结构。workload 一变，Owner 不变，选择改变。
\end{examplebox}

\section{比较表与边界}
\begin{center}\small
\begin{tabularx}{\linewidth}{L{2.4cm}YY}\toprule
策略 & 擅长 & 主要付费\\\midrule
折半/有序数组 & 静态查找、范围 & 插入移动、连续空间\\
BST/平衡树 & 动态有序查询 & 旋转、节点空间\\
Hash & 精确等值 & 冲突、装填、重散列、无序\\
B/B+ & 大规模块索引、范围 & 分裂合并与 I/O\\\bottomrule\end{tabularx}\end{center}

\section{做题控制协议}
先写至少两种候选，再逐项填查询、更新、空间、最坏与 I/O；任何“更快”都必须附带前提。若题目只给一个操作而没有 workload，保留多解并指出缺失信息，不擅自宣布唯一最佳结构。
\begin{corebox}{压缩复原}
四问：要查什么？多久更新？需不需要顺序？成本由时间、空间还是 I/O 主导？
\end{corebox}
\end{document}
